\documentclass[11pt,a4paper]{article}
\usepackage[T1]{fontenc}
\usepackage[utf8]{inputenc}
\usepackage{lmodern}
\usepackage[margin=26mm]{geometry}
\usepackage{amsmath,amssymb,amsthm}
\usepackage{array,booktabs,longtable,tabularx}
\usepackage{enumitem}
\usepackage[hidelinks]{hyperref}
\usepackage{microtype}
\setlength{\parindent}{0pt}
\setlength{\parskip}{5pt plus 1pt}
\setlist[itemize]{leftmargin=1.5em,itemsep=2pt,topsep=3pt}
\setlist[enumerate]{leftmargin=1.8em,itemsep=2pt,topsep=3pt}
\newcommand{\Cat}{\mathbf{Cat}}
\newcommand{\Set}{\mathbf{Set}}
\newcommand{\FinSet}{\mathbf{FinSet}}
\newcommand{\Ctx}{\operatorname{Ctx}}
\newcommand{\Cnt}{\operatorname{Cnt}}
\newcommand{\Kl}{\operatorname{Kl}}
\newcommand{\id}{\mathrm{id}}
\newcommand{\Hom}{\operatorname{Hom}}
\newcommand{\im}{\operatorname{im}}
\newcommand{\HistProb}{\mathbf{HistProb}}
\newcommand{\FinStoch}{\mathbf{FinStoch}}
\newcommand{\Dist}{\operatorname{Dist}}
\newcommand{\op}{\mathrm{op}}
\newcommand{\E}{\mathcal E}
\newcommand{\D}{\mathcal D}
\newcommand{\C}{\mathcal C}
\newcommand{\K}{\mathcal K}
\newcommand{\Lcat}{\mathcal L}
\newcommand{\SourceStatus}[1]{\textbf{#1}}
\title{Generalizing contextads\\[4pt]\large Scoping and strongest source audit}
\author{}
\date{For review\quad 1 November 2026}

\begin{document}
\maketitle
\vspace{-1.5em}

\textbf{Scope and status.} This is the authorized first substantive obligation
of the provisional research contract: a targeted strongest-source audit and a
typed scoping proposal. It does not adopt a theorem package, prove the proposed
main theorem, or discharge the remaining six obligations. The existing contract
has not been changed. The audit and candidate formulation are delivered for
review; the first gate is not thereby accepted.

\textbf{Main finding.} Substantial parts of the initially attractive target are
already established: dependent context composition, mixed opwreath and
two-sided bicategorical composition, law--lifting equivalences,
factorization-derived laws and double-category recognition, and several
dependent-effect obstructions. A newly located telescope
manuscript supplies a narrower positive comparison under a comprehension
\emph{equivalence}; it does not settle general PRA transposition or mixed
context/content interaction. The proposed mixed arrows can be typed precisely,
but an independently checkable interaction criterion at the desired retained-cell
level is not yet established. Even the tempting finite-set image-factorization
recipe fails associativity.

\textbf{Recommendation.} Do not freeze the generic claim that a suitably
factorizing contextad and contentad compose. Review the typed dependent-crossing
branch below, with its explicit open obligations, against the stronger sources.
The potentially substantive target is a \emph{specific dependent representation
and comparison theorem}, including a noncircular entry test and the exact indexed structure and
preservation requirements. Merely retaining cells or giving bounded recognition
is not enough: existing theorems already do that for substantial classes.
Whether the particular increment is new, already implied by existing
theorems, or unavailable for the intended examples remains a question, not a
conclusion of this audit.

\section{Programme and boundaries}

The programme is to generalize contextads through theoretical results and diverse
applications. Robots with discovered semantics and partial recovery are one stress
test, not the programme's defining object. Adaptive probabilistic computation,
typed protocols with allocated capabilities, and eventually differentiable or
noncartesian settings remain relevant. Later expansion is allowed through explicit
scope changes rather than being silently included in this first workstream.

The first provisional workstream remains dependent mixed context/content
composition with limited recognition. It should cover a parameterized class,
not merely a finite example. At least one later application must make the next
requirement vary with an actual returned result, and at least one must be
non-robotic and require no recovery. Existing one-sided constructions and the
ordinary fixed-comonad/fixed-monad control must be recovered at declared equality
levels. These are desired contents, not established achievements.

Three meanings of dependence must not be conflated:
\begin{enumerate}
\item A command changes a declared pre/post-state index. Parameterised effects
already support this extensively.
\item A realized result selects a later fibre or admissible requirement, with a
dependent sum or joint execution witness retained. Ordinary contextads already
support significant instances of this.
\item An arbitrary dependent section is extended across the monad unit
\(A\to TA\), where its type can inspect the whole effectful value. This is a
stronger dependent Kleisli-extension demand and can fail even when forward
sampling with result-indexed requirements works.
\end{enumerate}
The proposed first branch concerns (2), while comparing (1) and avoiding an
unannounced demand for (3).

\section{Evidence policy and the strongest baselines}

\subsection{What was checked}

The audit inspected the primary definitions and the indicated proof arguments,
not only abstracts. The companion source manifest records titles, versions,
public URLs and hashes of retrieved primary PDFs. The package contains no
primary-source PDFs or extracted copies. The inspection is targeted, not a
full-source dependency certification, machine formalization, or exhaustive
prior-art search. A statement that no mixed theorem was found in the inspected
pass is not a claim that none exists elsewhere.

The ledger uses four applicability classes: \textbf{A}, already proved and
applicable at its stated level; \textbf{H}, usable only after additional
hypotheses or an explicit bridge; \textbf{S}, stated but the needed proof is not
verified in the inspected source; and \textbf{U}, unresolved in this audit.
Proof status and applicability are separate: a sound established theorem may be
class H for this proposed construction. A definition is inspected data, not a
newly proved theorem.

\subsection{Contextads and contentads}

Capucci--Myers Definitions 3.15, 3.26 and 3.17 already give indexed context
requirements, their dependent accumulation and contextful composition. The
extension maps \(\varepsilon\) and \(\delta\) need not be invertible; the grade
unitors/associator and specified transport comparisons do. Example 3.20 uses a
comprehension category with unit, strong dependent sums and Beck--Chevalley;
Example 3.21 gives display-span examples. Thus dependence, dependent sums, partial
maps, or noninvertible extension structure alone are not new \cite{CM}.

Theorem 5.16 is a structure-lifting result for a colaxly or strongly
\(T\)-structured contextad under Definition 5.13, not an automatic interaction
with any second effect. Section 7, Definitions 7.9--7.10, already allows coherent
transport without universal cartesian lifting, but retains invertible transport,
unit-context and combination comparisons. Section 8 explicitly proposes mixed
contextad/contentad laws and factorization-based recognition. These proposals
must be attributed rather than rediscovered \cite{CM}.

Definitions 6.1--6.3 use right-opfibrant spans and the EM-side presentation for
contentads, rather than the contextad's left-fibrant/KL-side presentation. The
input pair is therefore not automatically two endopseudomonads on one already
specified 2-category. A common ambient bridge, with exact opposite conventions,
is part of the work.

There is also a proof-status boundary inside the source: Theorem 4.12 supplies an
explicit local 2-categorical isomorphism, whereas Section 4.4 and footnote 17
leave a stronger full KL tricocompletion identification conjectural. Remark
4.20 takes a higher wreath-product formula for granted by analogy. These caveats
limit imports of the stronger universal property; they do not invalidate the
explicit contextful composition formulas.

\subsection{Polynomial and PRA transposition}

Capucci--Myers Appendix A, Theorem A.4, states an identity-on-objects
polynomial Kleisli--\(\Ctx\) isomorphism. Remark A.5 explicitly leaves its proof
out and suggests a more general PRA route. The 2025 slides state the broader
PRA result using a left adjoint \(L:\C/T1\to\C\) and the shape/body
correspondence, but supply no complete proof. The inspected source status is
\textbf{S}, not ``false'' and not ``proved and freely importable''
\cite{CM,Slides}. A direct proof in a chosen class would be an attributed
verification or specialization unless an additional increment is identified.

Kamiya's October 2026 manuscript provides actual proofs of related, narrower
results \cite{Kamiya}. The relevant arguments in Sections 2--4 and Appendix
A.7--A.8 were inspected; no substantive gap was found in that focused reading,
without certifying every imported result. Its strict telescope-algebra result,
Corollary 2.20, concerns strict cartesian span maps with normal cleavage and
strict contextads with identity \(\varepsilon,\delta,\lambda,\rho,\alpha\).
Proposition 3.2 gives a monoidal Grothendieck extension for split left-fibrant
spans,
\[
  P_M(Y)(B)=\int_{P\in M_B}Y(B\mathbin{\odot}P).
\]
Proposition 3.5 and Theorem 4.2 assume a supplied display contextad and a
\emph{fibrewise comprehension equivalence} onto it, with the specified pullback,
composition and selected-square closure conditions. Transport and a thinness
lemma force the coherence; the extension maps are invertible in this subclass.
Proposition 3.6 compares the Grothendieck pseudomonad's Kleisli construction
\emph{on representables} with the contextful construction, identity on the
base/tight arrows. It does not compare the entire Kleisli bicategory on arbitrary
indexed categories, reconstruct an arbitrary computational double category,
prove arbitrary PRA transposition, or provide a mixed law.

The manuscript's declared AI-assisted provenance is not used as a mathematical
verdict. Its arguments are assessed on their content. A cited companion,
\emph{The contextad of double telescopes}, was not retrieved by the targeted
search and is not imported here.

\subsection{Mixed opwreaths are a stronger direct control}

Street's Definitions 3 and 5 distinguish a mixed wreath, a comonad in
\(\mathrm{EM}(\mathcal B)\), from a mixed \emph{opwreath}, a comonad in
\(\mathrm{KL}(\mathcal B)\), for a strict 2-category \(\mathcal B\). The latter
is the relevant variance for arrows \(GA\to TB\). For a monad \(T\), it supplies
\[
 \zeta:GT\Rightarrow TG,\qquad
 d:G\Rightarrow TGG,\qquad e:G\Rightarrow T,
\]
subject to seven axioms; \(G\) need not already be an ordinary comonad.
The construction and Proposition 7 give mixed Kleisli arrows and the representing
property. For \(f:GA\to TB\), \(g:GB\to TC\), the typed convolution is
\begin{align*}
 GA&\xrightarrow{d_A}TGGA\xrightarrow{TGf}TGTB
 \xrightarrow{T\zeta_B}TTGB\xrightarrow{TTg}TTTC
 \xrightarrow{\mu_C^{(3)}}TC,
\end{align*}
with identity \(e_A\). Proposition 6 relates opwreaths and wreaths by mates under
\(G\dashv S\); that mate calculation is left as a string-calculus exercise in
the source. Effectful duplication and discharge are therefore already covered
\cite{Street}. A theorem for this strict 2-categorical input does not automatically
specialize to the displayed-span tricategorical input; that bridge must be typed.

\subsection{A two sided bicategorical construction is also established}

Garner's Definition 14 already gives the two-sided Kleisli bicategory for a
pseudocomonad \(G\), pseudomonad \(T\) and mixed pseudo-distributive law on a
bicategory, with hom-categories \(\mathcal B(GA,TB)\) and the familiar five-map
composition. The generic associativity/unit constraints are not displayed in
full there. More importantly, Sections 3--4 actually derive a mixed law on
\(\mathrm{Mod}\) using double clubs and suitable matchings. Proposition 20 uses
pullback-created lifts; Proposition 38 and Corollary 39 verify all ten coherence
axioms by local-monomorphism cancellation. Those arguments were inspected
\cite{Garner}. The evaluation-at-one equivalence and monoidal structure invoke
results from the companion \emph{Double clubs}, not independently audited here.
Thus neither bicategorical mixed packaging nor the abstract idea of deriving a
law from extra structure is a new frontier. The remaining question is the
\emph{specific} identification and compatibility of the contextad/contentad
dependent data with such established inputs, including whether the club or
factorization hypotheses already imply the desired conclusion.

\subsection{Laws and liftings do not by themselves give an entry criterion}

Walker's Theorem 5.0.1 equates six presentations of a pseudodistributive law,
including compatible extension to a Kleisli bicategory. Theorem 4.1.8 needs
an additional compatibility W5, not merely a pseudomonad morphism and
opmorphism. Structural modifications are invertible; the law transformation
itself need not be. The inspected proof explicitly constructs the equivalences
but describes itself as a sketch and omits some modification checks
\cite{Walker}. A law or fully coherent lifting supplied as input is not an
independently derived criterion.

Gambino--Lobbia's Theorems 2.5 and 3.4 construct the Gray-categorical formal
lifting machinery; Theorem 4.5 identifies a pseudodistributive law with a
lifting to pseudoalgebras. Corollary 4.6 gives the right-module dual and, where
represented, a Kleisli extension. The proof reconstructs the law from free
pseudoalgebras; some routine coherence and quasi-inverse checks are omitted
\cite{GL}. These are strong established anti-reformulation tests. Strictifying
an ambient category alone does not prove preservation of the chosen displays,
pure arrows or cells.

\subsection{Factorization and pullback laws are already constructive baselines}

Rosebrugh--Wood Section 3.3 derives a crossing law by refactorizing composites
under \emph{strict unique} factorization through wide subcategories; Theorem
3.8 reconstructs the original category by an identity-on-objects isomorphism.
Sections 5.7--5.9 handle classical chosen factorizations relative to isomorphisms.
The inspected proofs derive equations from uniqueness. Strict factorization is
not the same definition as an ordinary orthogonal factorization system
\cite{RW}. Consequently, ``coherence follows from unique normal forms'' is not
itself the proposed advance.

Lack's Theorem 4.6 provides PROP recognition with factorization unique modulo
permutations, under wide, symmetry-containing, tensor-closed subcategories.
Sections 5.3 and 5.4 construct pullback and pushout laws by pasting. These results
work at the stated isomorphism-class level, discarding span/cospan 2-cells
\cite{Lack}. Their hypotheses and proofs were inspected, not all later algebraic
presentation examples. A retained-cell or dependent-context comparison may still
be substantive, but its novelty cannot be inferred from those papers' limits.

Fong--Zanasi's extended version supplies a further positive control. Assumption
1 requires pullbacks and pushouts, a costable \((\E,\mathcal M)\) factorization
with \(\mathcal M\) monic, and a pullback-stable subcategory containing
\(\mathcal M\) with the specified pushout-of-pullback comparison in
\(\mathcal M\). Theorem 3.1 presents corelations as a categorical pushout;
Appendix A's proof was inspected. Corollary 3.4 verifies the abelian case and
identifies corelations with relations; Section 4.7 covers linear relations
\cite{FZ}. A vector-space positive model cannot count relational composition
as new. The original Mal'cev/difunctionality characterization remains a
comparison lead whose primary proof was not inspected here.

\subsection{Double category factorization recognition is not a generic gap}

\v{S}t\v{e}p\'an's thesis Theorem 3.3.7 identifies codomain-discrete double
categories with strict factorization systems. Here each top-right corner has a
unique square filler. Theorem 3.3.33 identifies factorization double categories
with orthogonal factorization systems. Definition 3.3.24 requires invariance,
top-right corner fillers with every square bicartesian, and joint monicity of
bottom-left corners, implying flatness. The proofs reconstruct squares and
natural isomorphisms of double functors, not only objects; both theorem proofs
were inspected \cite{Stepan}. Consequently, a claimed increment must name the
additional indexed context/content structure or a precise non-flat comparison
outside those hypotheses. Passing from categories to flat doubles, retaining
squares, or making recognition functorial is already a substantial established
baseline. The thesis gives accessible proof text of the relevant work even
though the journal PDF retrieval failed in this pass.

\subsection{Parameterised and dependent effects}

Atkey's parameterised monad has
\(T:S^{\op}\times S\times\C\to\C\), dinatural unit/multiplication and strength
on a finite-product base. Propositions 1--3 and Theorem 1 construct and
characterize parameterised Kleisli/Freyd structure, with closure requiring
Kleisli exponentials. Theorem 3 supplies inverse lifting constructions.
The construction/proof paragraphs were inspected; routine checks are abbreviated,
and the referred syntactic-completeness proof was not re-audited
\cite{Atkey}. State-type changes and framed resources are already standard;
result-indexed families need a precise encoding comparison, not an absolute
claim that parameterised effects cannot represent them.

V\'ak\'ar's Theorem 5.2.9 constructs a dependent CBPV minus model from an indexed
monad using indexed Eilenberg--Moore algebras. Definition 5.3.2 adds dependent
Kleisli extensions with unit, composition and ordinary-extension compatibility;
Remark 5.3.4 derives dependent costrength. These are extra data. Sections
5.3.3.1--3 give explicit writer, reader and state non-models, and page 214 gives
a writer repair by a codiscrete setoid monoid, restricting what predicates can
distinguish. Those proofs and repair were inspected; the whole syntactic
completeness chain was not \cite{Vakar}. Neither unrestricted dependent
extension nor observation-restricting repair is an untouched research idea.

\subsection{Laxity is a separate branch}

Lack--Street Sections 4--5 already construct skew Kleisli bicategories.
Proposition 5.2 derives redundant axioms only under invertibility assumptions.
Theorem 8.1 normalizes a skew monoidal category using right-unit modules and
specified reflexive coequalizers preserved by right tensoring; Theorem 8.2
adds right-adjoint/opmonadic and hom-preservation hypotheses. The normalization
construction and proof were inspected \cite{LS}. Replacing pseudo comparisons
by noninvertible ones therefore requires a separate theorem and comparison,
not deletion of the word ``invertible'' from existing proofs.

\subsection{Compact applicability ledger}

\small
\begin{longtable}{@{}p{0.20\linewidth}p{0.10\linewidth}p{0.62\linewidth}@{}}
\toprule
Result or input & Class & Exact use and boundary \\
\midrule
\endhead
CM 3.26 and 3.17 & A & Existing dependent contextad and contextful composition; grade coherence is supplied. \\
CM 5.16 and 7.9 & H & Require the declared structured contextad or invertible transport comparisons; not arbitrary mixed interaction. \\
CM 4.4 and A.4 & S & Strong universal completion and polynomial transposition have explicit source proof caveats. \\
PRA slides & S & Precisely stated claimed general result; no complete proof in inspected slides. \\
Kamiya 2.20 & A/H & Strict telescope-algebra comparison only under the specified strictness and normal cleavage. \\
Kamiya 3.5, 3.6, 4.2 & H & Display contextad plus comprehension equivalence; Kleisli comparison only on representables. \\
Street 5 and 7 & A/H & Mixed opwreath construction already covers effectful duplication; higher displayed-span specialization is not supplied. \\
Garner 14, 38--39 & A/H & Two-sided Kleisli bicategory and a derived mixed law; the contextad-specific identification remains to map. \\
Walker 5.0.1 and GL 4.5 & A/H & Law--lifting equivalences are inherited machinery, not existence criteria for a missing law. \\
RW 3.8 and 5.9 & A/H & Strict or isomorphism-relative factorization laws; the dependent retained-cell bridge remains to check. \\
Lack 4.6 and 5.3--5.4 & A/H & Permutation-relative recognition and PB/PO laws at their quotient level. \\
\v{S}t\v{e}p\'an 3.3.7 and 3.3.33 & A/H & Square-preserving recognition for codomain-discrete and factorization doubles, with strong corner/flatness hypotheses. \\
FZ 3.1 and 3.4 & A/H & Explicit exactness assumptions; abelian relation/corelation model is an existing control. \\
Atkey 1--3 & A/H & Parameterised state/strength/closure constructions; no automatic mixed-contextad comparison. \\
V\'ak\'ar 5.2.9 and 5.3.2 & A/H & Indexed monad model and additional dependent extension; distinguish realized values from whole effects. \\
Lack--Street 8.1--8.2 & H & Skew normalization only with its exactness/preservation assumptions. \\
General dependent mixed criterion & U & No noncircular sufficient criterion, retained-cell bridge and bounded converse established for the intended common class. \\
\bottomrule
\end{longtable}
\normalsize

\section{Typed candidate charter}
\label{sec:charter}

\textbf{Proposal only.} The following fixes enough notation to expose the
interaction problem and compare sources. It is not an adopted definition of a
new generalized contextad, an accepted success criterion, or a proved construction.

\subsection{Base and one-sided inputs}

Take an ordinary category \(\C\) as the base, with the indexed presentations in
\(\Cat\). Size/universe choices must be fixed before a family theorem. Let
\[
 \K:\C^{\op}\longrightarrow\Cat,
 \qquad \Lcat:\C\longrightarrow\Cat
\]
be pseudofunctors, equipped respectively with a contextad and a contentad
structure. Chosen cleavages provide \(u^*\) and \(u_!\); a split presentation
may simplify calculation but is not a theorem that the original data are
literally equal to their split replacement.

Write \(e_A P\) for the context extension of \(P\in\K_A\), and \(d_A Q\)
for the content extension of \(Q\in\Lcat_A\). For \(u:A\to B\), the action
on the chosen cartesian/opcartesian lifts supplies
\[
 \bar u_P:e_A(u^*P)\longrightarrow e_B P,
 \qquad
 \widehat u_Q:d_AQ\longrightarrow d_B(u_!Q).
\]
There is no implicit pushforward on \(\K\) or pullback on \(\Lcat\).

The ordered context operations are
\[
 I_A\in\K_A,\quad P\otimes R\in\K_A\ (R\in\K_{e_AP}),\quad
 \varepsilon_A:e_AI_A\to A,\quad
 \delta_{P,R}:e_A(P\otimes R)\to e_{e_AP}R.
\]
The ordered content operations are
\[
 J_A\in\Lcat_A,\quad Q\diamond S\in\Lcat_A\ (S\in\Lcat_{d_AQ}),\quad
 \eta_A:A\to d_AJ_A,\quad
 \mu_{Q,S}:d_{d_AQ}S\to d_A(Q\diamond S).
\]
This convention fixes the indexing of content accumulation rather than copying
the ambiguous tensor order in the prose around CM equation (6.1.8).
All grade coherence and cartesian/opcartesian preservation required of these
one-sided inputs is supplied data at this stage, not derived by the mixed proposal.

\subsection{Mixed arrows and the missing crossing}

The proposed mixed arrow \(A\rightsquigarrow B\) is
\[
 (P,Q,f),\qquad P\in\K_A,\quad Q\in\Lcat_B,\quad f:e_AP\longrightarrow d_BQ.
\]
For a following arrow \((R,S,g):B\rightsquigarrow C\), one has
\(R\in\K_B\), \(S\in\Lcat_C\), \(g:e_BR\to d_CS\).
The expression \(f^*R\) is \emph{ill-typed}: \(f\) lands in \(d_BQ\), not
\(B\). This is not corrected by calling the expression dependent.

A crossing must at least give, for each \(Q\in\Lcat_B\), \(R\in\K_B\),
\[
 \widehat R\in\K_{d_BQ},\qquad
 \widehat Q\in\Lcat_{e_BR},\qquad
 \chi_{Q,R}:e_{d_BQ}\widehat R\longrightarrow d_{e_BR}\widehat Q.
 \tag{X}
\]
Then set
\[
 P'=P\otimes f^*\widehat R,
 \qquad S'=S\diamond g_!\widehat Q.
\]
The proposed composite body is now genuinely typed:
\begin{align*}
 e_A P'
 &\xrightarrow{\delta}e_{e_AP}(f^*\widehat R)
 \xrightarrow{\bar f}e_{d_BQ}\widehat R
 \xrightarrow{\chi_{Q,R}}d_{e_BR}\widehat Q\\
 &\xrightarrow{\widehat g}d_{d_CS}(g_!\widehat Q)
 \xrightarrow{\mu}d_C S'.
 \tag{C}
\end{align*}
The candidate identity is \((I_A,J_A,\eta_A\varepsilon_A)\).
Writing (C) proves only that these arrows compose as a \emph{formula} once
(X) is given. It does not prove functoriality, associativity, units or preservation.

The crossing needs functoriality in context/content morphisms, compatibility
with change of base, both units, both dependent accumulations, and the resulting
higher compatibility of those comparisons. If these are simply required as
axioms, they are supplied mixed-law data. They do not pass the contract's
anti-tautology test. An independently given universal property may derive them,
but its exact hypotheses must be stated and checked against the factorization
and law-lifting baselines.

\subsection{Cells and maps of instances}

A cell structure is not harmless decoration. One readily typed candidate for
fixed endpoints has a pair \(\alpha:P\to P'\), \(\beta:Q\to Q'\) satisfying
\[
 d_B(\beta)\, f=f'\,e_A(\alpha).
 \tag{H}
\]
For squares over \(u:A\to A'\), \(v:B\to B'\), the analogous direction is
\(\alpha:P\to u^*P'\), \(\beta:v_!Q\to Q'\), satisfying
\[
 d_{B'}(\beta)\,\widehat v_Q\,f
 = f'\,\bar u_{P'}\,e_A(\alpha).
\]
These directions agree with CM Definition 3.17, equations (3.3.15)--(3.3.16),
and the contentful direction in Example 6.4. Vertical composition is typed;
horizontal composition and the one-sided recovery equivalences are still proof
obligations. Reversing a cell direction would change the comparison theorem.
For a pure base map \(u:A\to B\), the candidate horizontal representative is
\((I_A,J_B,\eta_Bu\varepsilon_A)\), while the tight category can remain literally
\(\C\). The horizontal pure representation need not be faithful for arbitrary
inputs; faithfulness is not assumed without a proof.

A map of instances must comprise a base functor together with indexed functors
on \(\K,\Lcat\), lineator/colineator comparisons for \(e,d\), comparisons for
their units and accumulations, and compatibility with (X). The first candidate
should use structure-preserving pseudo maps; allowing arbitrary lax maps is a
separate choice. No naturality theorem for these maps has been established here.

\subsection{Continuations and observation interfaces}

The proposed unrestricted semantic continuation class consists of the mixed
arrows just typed, with \(R\) genuinely living over the returned base \(B\).
An application may instead require an explicitly declared subcategory of
observation-adapted or resource-legal continuations. Such a restriction must be
closed under the eventual identities/composition and preserved by comparison;
it is not automatically licensed by retaining a latent history in the semantics.

If \(q:H\to O\) forgets history, a deterministic observable
\(r:H\to D\) is allowed through \(O\) only when it factors through \(q\).
For varying resource families, a family and its permitted structure must descend
along \(q\) with specified isomorphisms and compatible policies. Retention of a
proof witness does not give a controller access to it, physical possession of a
resource, permission, conditioning, resetting, or an undo operation.

\section{The factorization branch and its precise failure}

\subsection{An independently described source of crossing data}

A tractable candidate specializes \(\K_A\) to display maps \(p:X\to A\)
from a class \(\D\), and \(\Lcat_B\) to codisplay maps \(q:B\to Y\) from a
class \(\E\). Context reindexing is pullback; content transport is pushout;
accumulation is composition, with identities included and closure/preservation
hypotheses explicitly required.

A mixed arrow is then a diagram
\[
 A\xleftarrow{p\in\D}X\xrightarrow{f}Y\xleftarrow{q\in\E}B.
\]
For a next display \(r:Z\to B\), factor
\[
 q r=r'c',\qquad Z\xrightarrow{c'\in\E}W\xrightarrow{r'\in\D}Y.
 \tag{F}
\]
This supplies (X) with \(\widehat R=r'\), \(\widehat Q=c'\), and
\(\chi=\id_W\). Pull \(r'\) back along \(f\), and push \(c'\) out along
\(g:Z\to V\), to get the candidate new normal form
\[
 A\longleftarrow X\times_Y W\longrightarrow W+_ZV\longleftarrow C.
 \tag{N}
\]
The outer legs use the original displays/codisplays and their composition.
This description uses independently meaningful operations, but their separate
existence is not enough. One needs compatibility of factorization with the
pullback/pushout interchanges and repeated normalization, including retained
maps of presentations. Neither (F) nor ordinary uniqueness up to isomorphism
already establishes those particular diagrams.

\subsection{A finite exact countercheck}

The following calculation rejects a plausible overstrong scoping assumption,
not all mixed constructions and not the already-established factorization-law
theorems. Let \(\C=\FinSet\), displays be inclusions and codisplays be
surjections; use image factorization in (F), ordinary pullbacks and pushouts in
(N). Present an arrow by a partial domain \(X\subseteq A\), a quotient
\(q:B\twoheadrightarrow Y\), and \(f:X\to Y\).

Let
\[
 A=\{a_0,a_1\},\quad B=\{b_0,b_1,b_2,b_3\},\quad
 C=\{c_0,c_1,c_2\},\quad D=\{*\}.
\]
The first arrow is total, with quotient blocks
\(\{b_0,b_1\}\mapsto y_0\), \(\{b_2,b_3\}\mapsto y_1\), and
\(a_i\mapsto y_i\). The second is the total pure function
\[
 b_0\mapsto c_0,\quad b_1\mapsto c_1,\quad
 b_2\mapsto c_1,\quad b_3\mapsto c_2.
\]
The third is defined only at \(c_2\), sending it to \(*\).

Combining the first two by (N), the pushout identifies
\(y_0\sim c_0\sim c_1\sim y_1\sim c_2\); hence their quotient target is a
singleton and both inputs remain in the domain. Composing with the third arrow
then still leaves both \(a_0,a_1\) defined. In the opposite grouping, the
second and third arrows first restrict to \(\{b_3\}\); its image in the first
quotient is \(\{y_1\}\), so the final domain is only \(\{a_1\}\).
Thus
\[
 \operatorname{dom}(h\star(g\star f))=\{a_0,a_1\}
 \ne\{a_1\}=\operatorname{dom}((h\star g)\star f).
\]
The domains differ even up to isomorphism over \(A\). The accompanying script
uses finite sets and exact equivalence-class pushouts and checks both results;
this displayed finite argument, not numerical sampling, explains the failure.

The calculation shows that individual factorization, pullback and pushout
operations do not imply the required mixed interchange. It is not a new
published impossibility claim, does not rule out all one-sided encodings, and
has no demonstrated positive repair in the proposed application class. Therefore
it does not complete the later contracted obstruction-and-repair obligation.

\subsection{Positive candidates remain comparisons rather than conclusions}

Abelian or vector-space normal forms are plausible controls because linear
relations are composition-closed and the relation/corelation comparison is
already established under precise hypotheses \cite{FZ}. Any use here must
supply the map from (N), explain which witnesses/cells survive, and compare with
that known category. A nontrivial retained-cell or indexed-instance result, if
available, would be the proposed increment. It is not yet proved, and linear
closure alone would not satisfy the contrasting application requirements.

\subsection{One explicit presentation distinction}

A quotient to underlying relations can lose content distinctions even when the
relation does not change. Let \(A=\{a\}\), \(B=\{0,1,2\}\), with full domain
and the output of \(a\) equal to the class of \(0\). The target partitions
\(\{0\}\mid\{1\}\mid\{2\}\) and \(\{0\}\mid\{1,2\}\) yield the same
relation \(\{(a,0)\}\), but the presented data retain a noninvertible quotient-
coarsening square from the first to the second; no reverse quotient map exists.
This specifies exactly one lost distinction. It does not show that retaining it
is desirable or novel, nor that a relational equipment cannot retain it.
The comparison must be made against the square-preserving recognition results
just audited, not against a deliberately stripped-down relation category alone.

\section{Assumption and conclusion ledger for the proposed target}

Every prospective hypothesis must have a role. The following table records the
current role rather than asserting that all entries suffice.

\small
\begin{longtable}{@{}p{0.24\linewidth}p{0.63\linewidth}@{}}
\toprule
Proposed input & Role and present status \\
\midrule
\endhead
Common base and universe & Typing. Must specify the class of bases, size and allowed maps; an arbitrary common higher ambient has not been constructed. \\
\(\K,\Lcat,e,d\) & Typing and inherited one-sided laws. Chosen pseudoindexed data and action functors; no missing opposite-variance transport inferred. \\
One-sided units and accumulation & Supplied data with their native coherence and preservation requirements; not a mixed theorem's derived output. \\
Chosen pullbacks/pushouts & Structural data for display/codisplay candidate; closure of selected classes needed. Existence alone fails by the finite check. \\
Chosen factorization (F) & Candidate entry data. Need its exact functoriality/uniqueness level; ordinary image factorization is insufficient for (N). \\
Mixed PB/PO/factorization compatibility & Essential open proof obligation. Must be an independently checkable property, not the desired associativity repackaged. \\
Crossing (X) and law comparisons & If assumed, supplied mixed-law data and a formal baseline. To meet the research goal, a substantive part must be derived from independent inputs. \\
Comprehension equivalence & Optional strong imported hypothesis for a telescope/display branch; gives representation data rather than intrinsic reconstruction. \\
Strength or Beck--Chevalley & Must be specified where used, including exact direction and preservation property. Not implied by the word ``indexed''. \\
Cells and maps of instances & Proposed pseudo comparison level. Cell variance is typed and matches the one-sided sources; horizontal/base-change stability and recovery equivalences remain to prove. \\
Continuation interface & Application data. Needs identity/composition closure and descent/preservation proof; cannot be changed silently to make a quotient work. \\
Quotient or normal-form choice & Changes equality/presentation information. Record faithfulness, reflection or equivalence separately, and any discarded witnesses. \\
Lax structural comparisons & Not selected in the first candidate. Would require a separate skew/coherence target and revised comparisons. \\
\bottomrule
\end{longtable}
\normalsize

\textbf{Derived so far.} The displayed composite is well-typed given its
crossing and action maps; the finite image-factorization recipe has the stated
associativity counterexample; the distinctions among the source hypotheses
are recorded. No coherence theorem for the proposed mixed class follows from
these facts. In particular, no universal compatibility criterion, sharp positive
repair, naturality theorem or bounded recognition theorem has been proved here.

\section{Conservative controls and equality levels}

\subsection{One-sided controls}

With content trivialized, aim for a structure-preserving equivalence over the
base with \(\Ctx\); with context trivialized, aim for the analogous comparison
with \(\Cnt\). These must include cells, pure/tight arrows, identities and
composition. Literal equality is possible only after specified presentation
choices and proof; it is not the default. A calculation on arrow objects alone
does not establish either comparison at the double-category level.

\subsection{The ordinary fixed-functor control}

For an ordinary comonad \((G,\varepsilon,\delta)\), monad
\((T,\eta,\mu)\), and law \(\lambda:GT\Rightarrow TG\), the four laws are
\begin{align*}
 \lambda\,G\eta&=\eta G,& T\varepsilon\,\lambda&=\varepsilon T,\\
 \lambda\,G\mu&=\mu G\,T\lambda\,\lambda T,&
 T\delta\,\lambda&=\lambda G\,G\lambda\,\delta T.
\end{align*}
The fixed control has arrows \(f:GA\to TB\), identity
\(\eta_A\varepsilon_A\), and
\[
 g\star f=\mu_C\,Tg\,\lambda_B\,Gf\,\delta_A.
\]
A specialized mixed theorem must recover this displayed formula, not only an
unspecified equivalent category. Street's opwreath gives a stronger existing
control by setting \(d=\eta_{GG}\delta\), \(e=\eta\varepsilon\) in the ordinary
case. The pseudo analogue carries additional coherence; it is not inferred
solely from this ordinary calculation.

\subsection{History and law are different comparison levels}

The scope-corrected finite project report is an existing control, not a newly
established mixed theorem \cite{Finite}. Its proof notes were inspected for the
following narrow claims. Finite full-support probability families over
\(\FinSet\), with measure-preserving maps, form an ordinary contextad by
weighted dependent pairs. The horizontal bicategory \(\HistProb\) retains sample
histories; evaluation \(E:\HistProb\to\FinStoch\) gives the output law.
For each hom-category, the support sampler \(J\) is fully faithful and
\(E_{A,B}\dashv J_{A,B}\), with \(EJ=\id\). Globally, \(J\)'s comparison
\[
 J(g)\circ J(f)\longrightarrow J(g\circ f)
\]
forgets the intermediate sample and is generally noninvertible. Thus local
reflection is not global history-preserving equivalence. The hom-category
connected-component quotient recovers laws, not raw execution witnesses.

The same notes prove a narrowly represented deterministic-writer obstruction
with extensional normalized probability and specified pure-writer and grade-zero
stochastic inclusions; random grades or retained sample presentations change that
representation and repair the example. The tested mixture transport also fails
Section 7's invertible compositor and unit-context comparison. None of these
results forbids stochastic contextads in general.

Any future quotient must state exactly what continuations it preserves.
For a surjective finite history-forgetting map, deterministic observables descend
exactly when constant on its fibres. Unstructured finite resource families descend
up to fibrewise bijection exactly when their cardinalities are constant on each
fibre; labels, permissions, probabilities and policies require their own compatible
structure. Equal final success probabilities alone are weaker than such descent.

\section{Application probes and excluded shortcuts}

\begin{enumerate}
\item \textbf{Adaptive probabilistic computation.} Require a declared realized
outcome \(b\), a family of next requirements over \(b\), and an account of shared
latent information versus fresh conditional samples. The ordinary probability-family
contextad is a positive control, not by itself the mixed result. A probability
law on each resource fibre is data, not supplied by a bare category-valued fibre.
Empty positive-probability fibres need a stated restriction, abort or subprobability
semantics. No disintegration theorem is assumed.
\item \textbf{Typed protocols with allocation.} Require a formal world/capability
base, admissible maps, a returned-value-indexed requirement and the creation or
transfer of capabilities. Mere pre/post-state index change is already an Atkey
control. The comparison must say whether a fixed index plus a sum/existential or
a larger state space preserves the same continuation interface. No concrete
protocol model has yet been shown to satisfy the proposed mixed criterion.
\item \textbf{Linear or abelian examples.} Useful as exact categorical controls
for a proposed positive factorization branch. They do not replace the requirement
for two contrasting applications or establish novelty by relational closure.
\item \textbf{Robot and recovery stress test.} Optional later instantiation.
Separate latent execution evidence, observable receipts and allowed policies.
Compensation is not inversion. Specify the recovery domain and residual relation;
semantic discovery alone does not validate refinement of an earlier model.
\end{enumerate}

For every selected application, the next charter revision must provide an
explicit row identifying which of the three dependence notions is used, the
actual fibres and maps, the interaction hypothesis, the allowed continuations,
and the equality level. At least one nonidentity map of instances must exercise
the dependence/interaction in the eventual naturality theorem. None of these
fixtures is certified by this audit.

As an additional diagnostic, unrestricted dependent Kleisli extension for the
normalized finite-distribution monad on \(\Set\) cannot hold for all
\(\mathrm{Fam}(\Set)\) families. Take \(A=\{0,1\}\) and a family over
\(\Dist(A)\) that is singleton at Dirac laws and empty elsewhere. Along
\(\eta:A\to\Dist(A)\), its free-computation family has a section. At the fair
mixture its fibre is \(\Dist(\varnothing)=\varnothing\), so no global extension
exists. This is an elementary deduction in the pattern of V\'ak\'ar's definition,
not a probability theorem quoted from that source. It does not obstruct
forward-sampling dependent sums or the existing finite contextad.

\section{Recommended decision and candidate theorem package}

\subsection{What can and cannot be frozen now}

\textbf{Can be reviewed now:} the broad programme, the realized-value dependence
requirement, the typed arrow and crossing interfaces, the preservation/equality
standards, the strongest verified controls, and the need for bounded recognition.

\textbf{Not ready to freeze:} a sufficient structural class for (C), the chosen cell-level
ambient comparison, a compatible positive family that includes
the intended applications, or a novelty claim. The finite-set failure makes this
more than a missing exposition detail. A generic ``stable factorization'' label
would hide precisely the obligation still at issue.

\subsection{Preferred candidate for review}

The recommended candidate is a dependent mixed \emph{normal-form and comparison}
workstream based on (X)--(N), subject to an explicit gate before theorem execution:
\begin{enumerate}
\item Select a nontrivial class of display/codisplay presentations with independently
stated interchange/exactness conditions and at least one compatible positive
model. Determine exactly which conditions are already sufficient by the audited
factorization, relation/corelation and law-lifting results.
\item Specify the retained cells and maps of instances; prove or identify the
bridge into an existing mixed construction. Any coherence imported from a source
must have its hypotheses mapped, rather than being presented as new.
\item State a bounded recognition claim only for the chosen normal-form and
preservation class. For example, it may reconstruct a crossing from composition
that preserves designated pure arrows and normal forms. If an existing
factorization theorem already supplies that converse, count it as inherited and
identify the further comparison or boundary statement needed.
\item Preserve the one-sided controls and ordinary fixed formula at their declared
levels, then check the two contrasting application families. If the selected
class cannot accommodate them without material changes, return to the user
with the old and proposed targets side by side.
\end{enumerate}
This is a recommended investigation target, not a claim that these steps are
feasible in the full initial generality. The gate should not be called accepted
until the positive class, comparison level and residual increment are agreed.

\subsection{Alternative scoped choices}

\textbf{Display and comprehension branch.} Use a supplied display contextad and
comprehension equivalence, benefiting from the telescope results. Scope cost:
strong representation assumptions and invertible-extension subclass; no intrinsic
recognition of arbitrary computational double categories. The mixed opposite-
variance interaction remains to be supplied or derived.

\textbf{Interface and descent branch.} Focus on a declared family of effect
presentations and allowed continuation interfaces, deriving when an observation
quotient preserves dependent resource use. Scope cost: the interface is explicit
input; the elementary finite factorization test and existing dependent-effect
repairs are controls, not the new theorem. Adopting this instead of the broader
mixed family theorem would require an explicit scope decision.

\textbf{Existing-calculus specialization.} Establish a careful contextad/contentad
presentation of a known relational, parameterised or opwreath calculus, with
applications. This may be useful exposition and comparison, but must be labelled
as such. It does not silently replace the agreed goal of a substantive criterion
and boundary result.

\textbf{Lax or noncartesian expansion.} Retain noninvertible transport or move to a
different ambient setting. This remains a later option; it would reopen coherence,
normalization and control comparisons and cannot be selected by notation alone.

\subsection{Open items and stopping condition}

The audit leaves the following exact unresolved items: the independently
sufficient mixed interchange property; a successful dependent positive class;
the specifically indexed cell and change-of-instance bridge; one-sided cell recovery;
the nontrivial residual theorem after applying the strongest baselines; and the
concrete application fixtures. The full PRA proof, the stronger universal
tricocompletion claims, the unretrieved double-telescope companion and the
original Mal'cev characterization are separate literature items, not hidden
imports needed for the displayed calculation.

The present deliverable completes the targeted source/scoping pass and presents
its findings, including the negative result. It does not certify an exhaustive
audit or accepted charter. The proposed next scoping step is a focused hypothesis-mapping test: select one
positive class, test its PB/PO/factorization interchange against the existing
club, relation/corelation and double-factorization theorems, and state the exact
residual theorem or explain that the construction is inherited. The next decision
is whether to pursue that candidate, select an alternative explicitly, or seek a
broader comparison before freezing any theorem statement.

\begin{thebibliography}{99}
\bibitem{CM} Matteo Capucci and David Jaz Myers.
\emph{Contextads as Wreaths; Kleisli, Para, and Span Constructions as Wreath Products}.
arXiv:2410.21889v1, 2024. \url{https://arxiv.org/abs/2410.21889}.
\bibitem{Slides} Matteo Capucci. \emph{Contextads abridged}.
PSSL110 slides, May 2025. \url{https://matteocapucci.eu/assets/2025/05/pssl-110.pdf}.
\bibitem{Kamiya} Kotaro Kamiya.
\emph{The Grothendieck construction between the telescope monad and contextads}.
Hexagon 2610.00302v1, 17 October 2026.
\url{https://hexagonmath.org/2610.00302v1}.
\bibitem{Street} Ross Street. \emph{Wreaths, mixed wreaths and twisted coactions}.
arXiv:1610.07281. \url{https://arxiv.org/abs/1610.07281}.
\bibitem{Garner} Richard Garner. \emph{Polycategories via pseudo-distributive laws}.
arXiv:math/0606735v2. \url{https://arxiv.org/abs/math/0606735v2}.
\bibitem{Walker} Charles Walker. \emph{Presentations of pseudodistributive laws}.
arXiv:2102.12468v4. \url{https://arxiv.org/abs/2102.12468v4}.
\bibitem{GL} Nicola Gambino and Gabriele Lobbia.
\emph{On the formal theory of pseudomonads and pseudodistributive laws}.
Theory and Applications of Categories 37 (2021), 14--56.
\url{https://www.tac.mta.ca/tac/volumes/37/2/37-02.pdf}.
Author copy inspected:
\url{https://eprints.whiterose.ac.uk/id/eprint/170302/1/pseudomonads-tac.pdf}.
\bibitem{RW} Robert Rosebrugh and R. J. Wood.
\emph{Distributive Laws and Factorization}.
Journal of Pure and Applied Algebra 175 (2002), 327--353.
\url{https://mta.ca/~rrosebru/articles/fac.pdf}.
\bibitem{Lack} Stephen Lack. \emph{Composing PROPs}.
Theory and Applications of Categories 13 (2004), 147--163.
Inspected journal PDF revised 2021.
\url{https://tac.mta.ca/tac/volumes/13/9/13-09.pdf}.
\bibitem{Stepan} Miloslav \v{S}t\v{e}p\'an.
\emph{Lax structures in 2-category theory}. Doctoral thesis, 2025, Chapter 3.
\url{https://arxiv.org/abs/2504.04467v1}.
Includes the relevant factorization-double-category results published in
Theory and Applications of Categories 41 (2024), 551--592.
\bibitem{FZ} Brendan Fong and Fabio Zanasi.
\emph{Universal Constructions for (Co)Relations: categories, monoidal categories, and props}.
Extended paper, arXiv:1710.03894.
\url{https://arxiv.org/abs/1710.03894}.
\bibitem{Atkey} Robert Atkey. \emph{Parameterised Notions of Computation}.
Journal of Functional Programming 19 (2009), 335--376.
\url{https://bentnib.org/paramnotions-jfp.pdf}.
\bibitem{Vakar} Matthijs V\'ak\'ar. \emph{In Search of Effectful Dependent Types}.
Doctoral thesis, 2017.
\url{https://www.cs.ox.ac.uk/people/aleks.kissinger/theses/vakar-thesis.pdf}.
Companion paper: \url{https://arxiv.org/abs/1603.04298}.
\bibitem{LS} Stephen Lack and Ross Street. \emph{On monads and warpings}.
Cahiers de Topologie et G\'eom\'etrie Diff\'erentielle Cat\'egoriques
55 (2014), 244--266.
\url{https://cahierstgdc.com/wp-content/uploads/2017/05/LackStreet_55-4.pdf}.
\bibitem{Finite} \emph{Generalizing contextads: A finite stochastic baseline and one recovery application}.
Project report, scope-corrected 1 November 2026, version 1.
Existing project control; its detailed composition, local reflection and descent
proof notes were inspected in this audit. Not a primary literature novelty claim.
\end{thebibliography}
\end{document}
